\setcounter{section}{0}
\part*{IV. BỘ BA CÂU PHÂN LOẠI TRONG 76 ĐỀ THI THỬ THPT QG  2015}
\addcontentsline{toc}{part}{IV. BỘ BA CÂU PHÂN LOẠI TRONG 76 ĐỀ THI THỬ THPT QG  2015}
\section{Đề minh hoạ THPT 2015} 
\noindent\debai{
Giải bất phương trình $\sqrt{x^2+x}+\sqrt{x-2}\ge \sqrt{3(x^2-2x-2)}$.
}
\\
\debai{
\bd cho tam giác $OAB$ có các đỉnh $A,B$ thuộc đường thẳng $\Delta: 4x+3y-12=0$ và $K(6;6)$ là tâm đường tròn bàng tiếp góc $O$. Gọi $C$ là điểm nằm trên $\Delta$ sao cho $AC=AO$ và các điểm $C,B$ nằm khác phía nhau so với $A$. Biết $C$ có hoành độ bằng $\frac{24}5$, tìm tọa độ $A,B$.
}
\\
\debai{
 Cho $x\in \mathbb{R}$. Tìm GTNN của
$$ P=\dfrac{\sqrt{3(2x^2+2x+1)}}{3}+\dfrac{1}{\sqrt{2x^2+(3-\sqrt 3)x+3}}+\dfrac{1}{\sqrt{2x^2+(3+\sqrt 3)x+3}}. $$
}

\section{Đề Sở GDĐT Phú Yên}
\noindent\debai{
\bd cho hình vuông $ABCD$ có $M,N$ lần lượt là trung điểm của $BC$, $CD$. Tìm tọa độ $B, M$ biết $N(0;-2)$, đường thẳng $AM$ có phương trình $x+2y-2=0$ và cạnh hình vuông bằng $4$. 
}

\noindent\debai{
Giải hệ phương trình $\hpt{27x^3+3x+(9y-7)\sqrt{6-9y}=0\\ \frac 13x^2+y^2+\sqrt{2-3x}-\frac{109}{81}=0}$.
}

\noindent\debai{
 Tìm GTLN và GTNN của biểu thức $P=5^{2x}+5^y$, biết $x\ge 0$, $y\ge 0, x+y=1$.
}

\section{THTT số 453 tháng 04 năm 2015}

\noindent\debai{
\bd hãy tính diện tích tam giác $ABC$ biết $H(5;5)$, $I(5;4)$ lần lượt là trực tâm và tâm đường tròn ngoại tiếp tam giác $ABC$ và cạnh $BC$ nằm trên đường thẳng $x+y-8=0$.
}

\noindent\debai{
 Giải phương trình $(x-\ln x)\sqrt{2x^2+2}=x+1$.
}

\noindent\debai{
Cho $0<x<y<z$. Tìm GTNN
$$ P=\dfrac{x^3z}{y^2(xz+y^2)}+\dfrac{y^4}{z^2(xz+y^2)}+\dfrac{z^3+15x^3}{x^2z}. $$
}

\section{THPT Số 3 Bảo Thắng}

\noindent\debai{
Giải hệ phương trình $\hpt{\sqrt{2x-y-1}+\sqrt{3y+1}=\sqrt x+\sqrt{x+2y}\\ x^3-3x+2=2y^3-y^2}.$
}

\noindent\debai{
\bd cho hình vuông $ABCD$ có tâm $I\left(\frac 72;\frac 32\right)$. Điểm $M(6;6)\in AB$ và $N(8;-2)\in BC$. Tìm tọa độ các đỉnh của hình vuông.
}

\noindent\debai{
Cho $x,y,z\in (0;1)$ thỏa mãn $(x^3+y^3)(x+y)=xy(1-x)(1-y)$. Tìm GTLN
$$ P=\dfrac{1}{\sqrt{1+x^2}}+\dfrac{1}{\sqrt{1+y^2}}+3xy-(x^2+y^2). $$
}

\section{THPT Bố Hạ (Bắc Giang)}

\noindent\debai{
\bd cho tam giác $ABC$ có trực tâm $H(-1;3)$, tâm đường tròn ngoại tiếp $I(-3;3)$, chân đường cao kẻ từ đỉnh $A$ là điểm $K(-1;1)$. Tìm tọa độ $ABC$.
}

\noindent\debai{
Giải hệ phương trình $\hpt{x^2(x-3)-y\sqrt{y+3}=-2\\ 3\sqrt{x-2}=\sqrt{y(y+8)}}$.
}

\noindent\debai{
Cho $x,y,z\in \mathbb{R}$ thỏa mãn $x^2+y^2+z^2=9$, $xyz\le 0$. CMR $2(x+y+z)-xyz\le 10$.
}

\section{THPT Chu Văn An (Hà Nội)}

\noindent\debai{
\bd cho tam giác $ABC$ có tâm đường tròn ngoại tiếp là $I(-2;1)$ và thỏa mãn điều kiện $\goc{AIB}=90^o$, chân đường cao kẻ từ $A$ đến $BC$ là $D(-1;-1)$, đường thẳng $AC$ đi qua $M(-1;4)$. Tìm tọa độ $A,B$ biết $A$ có hoành độ dương.
}

\noindent
\debai{
Giải BPT $3(x^2-2)+\dfrac{4\sqrt 2}{\sqrt{x^2-x+1}}>\sqrt x(\sqrt{x-1}+3\sqrt{x^2-1})$. 
}

\noindent
\debai{
Xét các số thực dương $x,y,z$ thỏa mãn điều kiện $2(x+y)+7z=xyz$. Tim GTNN $$S=2x+y+2z.$$
}

\section{THPT chuyên Hà Tĩnh}

\noindent
\debai{
\bd cho hình bình hành $ABCD$ có $N$ là trung điểm $CD$ và $BN$ có phương trình $13x-10y+13=0$; điểm $M(-1;2)$ thuộc đoạn thẳng $AC$ sao cho $AC=4AM$. Gọi $H$ là điểm đối xứng của $N$ qua $C$. Tìm tọa độ $A,B,C,D$ biết $3AC=2AB$ và $H\in \Delta: 2x-3y=0$.
}

\noindent
\debai{
Giải hệ phương trình $\hpt{x^2+(y^2-y-1)\sqrt{x^2+2}-y^3+y+2=0\\ \sqrt[3]{y^2-3}-\sqrt{xy^2-2x-2}+x=0}.$
}

\noindent
\debai{
Cho $a\in[1;2]$. CMR $(2^a+3^a+4^a)(6^a+8^a+12^a)<24^{a+1}$.
}

\section{THPT Đặng Thúc Hứa (Nghệ An)}

\noindent
\debai{
\bd cho hình bình hành $ABCD$ có $\goc{ABC}$ nhọn, $A(-2;-1)$. Gọi $H,K,E$ lần lượt là hình chiếu vuông góc của $A$ trên các đường thẳng $BC, BD, CD$. Đường tròn $(C):$ $x^2+y^2+x+4y+3=0$ ngoại tiếp tam giác $HKE$. Tìm tọa độ $B,C,D$ biết $H$ có hoành độ âm, $C$ có hoành độ dương và nằm trên đường thẳng $x-y-3=0$.
}

\noindent
\debai{
Giải hệ phương trình $\hpt{3\sqrt{y^3(2x-y)}+\sqrt{x^2(5y^2-4x^2)}=4y^2\\ \sqrt{2-x}+\sqrt{y+1}+2=x+y^2}$.
}

\noindent
\debai{
Cho $a,b,c>0$ thỏa mãn $4(a^3+b^3)+c^3=2(a+b+c)(ac+bc-2)$. Tìm GTLN
$$ P=\dfrac{2a^2}{3a^2+b^2+2a(c+2)}+\dfrac{b+c}{a+b+c+2}-\dfrac{(a+b)^2+c^2}{16}. $$
}

\section{THPT Đông Đậu (Vĩnh Phúc)}

\noindent
\debai{
Giải hệ phương trình $\hpt{\sqrt{x+2y+1}-2x=4(y-1)\\ x^2+4y^2+2xy=7}$.
}

\noindent
\debai{
\bd cho tam giác $ABC$ có phương trình đường thẳng $AB$: $2x+y-1=0$, phương trình $AC:$, $3x+4y+6=0$ và điểm $M(1;-3)$ nằm trên $BC$ thỏa mãn $3MB=2MC$. Tìm tọa độ trọng tâm $G$ của tam giác $ABC$.
}

\noindent
\debai{
Tìm tất cả các giá trị của $m$ để bất phương trình sau có nghiệm trên $[0;2]$
$$ \sqrt{(m+2)x+m}\ge |x-1|. $$
}

\section{THPT chuyên Hưng Yên}

\noindent
\debai{
\bd cho hình vuông $ABCD$ có $A(-1;2)$. Gọi $M, N$ lần lượt là trung điểm của $AD$ và $DC$; $K=BN\cap CM$. Viết phương trình đường tròn ngoại tiếp tam giác $BMK$, biết $BN$ có phương trình $2x+y-8=0$ và điểm $B$ có hoành độ lớn hơn $2$. 
}

\noindent
\debai{
Giải hệ phương trình $\hpt{(1-y)\sqrt{x^2+2y^2}=x+2y+3xy\\ \sqrt{y+1}+\sqrt{x^2+2y^2}=2y-x}$.
}

\noindent
\debai{
Cho $x,y,z>0$ thỏa mãn $5(x^2+y^2+z^2)=9(xy+2yz+zx)$. Tìm GTLN
$$ P=\dfrac{x}{y^2+z^2}-\dfrac{1}{(x+y+z)^3}. $$
}

\section{THPT chuyên Lê Hồng Phong (Hồ Chí Minh)}

\noindent
\debai{
\bd cho hình thang $ABCD$ có đáy lớn $CD=2AB$, điểm $C(-1;-1)$, trung điểm của $AD$ là $M(1;-2)$. Tìm tọa độ $B$, biết diện tích tam giác $BCD$ bằng $8$, $AB=4$ và $D$ có hoành độ dương. %(Đề gốc là nguyên dương, nhưng xét thấy không cần giả thiết nguyên nên bỏ).
}

\noindent
\debai{
Giải hệ phương trình $\hpt{2+9.3^{x^2-2y}=(2+9^{x^2-2y}).5^{2y-x^2+2}\\ 4^x+4=4x+4\sqrt{2y-2x+4}}.$
}

\noindent
\debai{
Cho $x,y,z>0$ thỏa mãn $x+y+1=z$. Tìm GTNN
$$ P=\dfrac{x}{x+yz}+\dfrac{y}{y+zx}+\dfrac{z^2+2}{z+xy}. $$
}

\section{THPT Lê Xoay}

\noindent
\debai{
\bd cho tam giác $ABC$. Đường phân giác trong góc $A$ có phương tình $d$: $x-y+2=0$, đường cao hạ từ $B$ có phương trình $d'$: $4x+3y-1=0$. Biết hình chiếu của $C$ lên $AB$ là điểm $H(-1;-1)$. Tìm tọa độ $B,C$.
}

\noindent
\debai{
Giải hệ phương trình $\hpt{xy(x+1)=x^3+y^2+x-y\\ 3y(2+\sqrt{9x^2+3})+(4y+2)(\sqrt{1+x+x^2}+1)=0}$.
}

\noindent
\debai{
Cho $a,b,c>0$ thỏa mãn $a+b+c=2$. Tìm GTLN
$$ S=\sqrt{\dfrac{ab}{ab+2c}}+\sqrt{\dfrac{bc}{bc+2a}}+\sqrt{\dfrac{ca}{ca+2b}}. $$
}

\section{THPT Lục Ngạn số 1 (Bắc Giang)}

\noindent
\debai{
Giải hệ phương trình $\hpt{y^2-x\sqrt{\dfrac{y^2+2}{x}}=2x-2\\ \sqrt{y^2+1}+\sqrt[3]{2x-1}=1}$.
}

\noindent
\debai{
 \bd cho $A(2;1)$, $B(-1;-3)$ và hai đường thẳng $d_1$: $x+y+3=0$, $d_2$: $x-5y-16=0$. Tìm tọa độ $C\in d_1$ và $D\in d_2$ sao cho $ABCD$ là hình bình hành.
}

\noindent
\debai{
Cho $x,y\in \mathbb{R}$ thỏa mãn $x^2+y^2+xy=3$. Tìm GTLN và GTNN $P=x^3+y^3-3x-3y$.
}

\section{THPT Lương Ngọc Quyến (Thái Nguyên)}

\noindent
\debai{
\bd cho hình vuông $ABCD$. $F\left(\dfrac{11}2;3\right)$ là trung điểm $AD$. $EK$: $19x-8y-18=0$ với $E$ là trung điểm $AB$, $K$ thuộc cạnh $CD$ sap cho $KD=3KC$. Tìm tọa độ $C$ biết $x_E<3$.
}

\noindent
\debai{
Giải hệ phương trình $\hpt{|x-2y|+1=\sqrt{x-3y}\\ x(x-4y+1)+y(4y-3)=5}$.
}

\noindent
\debai{
Cho $a,b,c>0$. CMR
$$ \dfrac{a^2+1}{4b^2}+\dfrac{b^2+1}{4c^2}+\dfrac{c^2+1}{4a^2}\ge \dfrac{1}{a+b}+\dfrac{1}{b+c}+\dfrac{1}{c+a}. $$
}

\section{THPT Lương Thế Vinh (Hà Nội)}

\noindent
\debai{
\bd cho hình thang $ABCD$ vuông tại $A,D$; diện tích hình thang bằng $6$; $CD=2AB$, $B(0;4)$. Biết $I(3;-1)$, $K(2;2)$ lần lượt nằm trên đường thẳng $AD$ và $DC$. Viết phương trình đường thẳng $AD$ biết $AD$ không song song với trục tọa độ.
}

\noindent
\debai{
Giải hệ phương trình $\hpt{x+\sqrt{x(x^2-3x+3)}=\sqrt[3]{y+2}+\sqrt{y+3}+1\\ 3\sqrt{x-1}-\sqrt{x^2-6x+6}=\sqrt[3]{y+2}+1}$.
}

\noindent
\debai{
Cho $x,y>0$ thỏa mãn $x-y+1\le 0$. Tìm GTLN
$$ T=\dfrac{x+3y^2}{\sqrt{x^2+y^4}}-\dfrac{2x+y^2}{5x+5y^2}. $$
}



